\documentclass[11pt]{article}

\usepackage[margin=1in]{geometry}
\usepackage{amsmath, amssymb, mathtools}
\usepackage{booktabs}
\usepackage{array}
\usepackage{longtable}
\usepackage{tabularx}
\usepackage{graphicx}
\usepackage{float}
\usepackage{siunitx}
\usepackage{enumitem}
\usepackage[hidelinks]{hyperref}

\sisetup{per-mode=symbol, group-separator={,}}
\setlength\LTleft{0pt}

\title{Stage 3 Fixed-Wing Tail Refinement and AeroSandbox Verification}
\author{AA146 Capstone Optimizer Notes}
\date{\today}

\begin{document}
\maketitle

\begin{abstract}
This document records the current Stage~3 fixed-wing tail refinement for the AA146 distributed-propulsion aircraft. Stage~3 freezes the Stage~1/2 propulsion, main-wing, airfoil, flap, and aileron selections. It sizes the stabilizer planform and tail arm for cruise trim, cruise drag, mass, static margin, and tail-volume constraints, then sizes the elevator and rudder chord fractions for slow-flight control authority. The present recommended configuration uses 10 propellers of diameter \SI{5.50}{in}, the DAE51 main airfoil, a horizontal tail span of \SI{0.606}{m}, and a vertical tail height of \SI{0.377}{m}. The aerodynamics remain tied to the locked gross flight mass of \SI{5.000}{kg}; the current component-estimated mass before ballast or payload is \SI{2.491}{kg}, with \SI{0.045}{kg} of added tail foam mass, cruise electrical power of \SI{112.1}{W}, and low-speed electrical power of \SI{188.4}{W}.
\end{abstract}

\section{Executive Summary}
Tables~\ref{tab:exec-configuration}, \ref{tab:exec-performance}, and \ref{tab:exec-margins} summarize the aircraft as it should be reviewed at the configuration level. Stage~1/2 selections are listed explicitly because Stage~3 treats them as frozen inputs; Stage~3 only refines the fixed-wing tail, trim, drag buildup, and reporting around that selected aircraft.

\begin{longtable}{@{}p{0.29\linewidth}p{0.30\linewidth}p{0.34\linewidth}@{}}
\caption{Executive aircraft configuration summary.}\label{tab:exec-configuration}\\
\toprule
Item & Current value & Source or design note \\
\midrule
\endfirsthead
\toprule
Item & Current value & Source or design note \\
\midrule
\endhead
Gross flight mass & \SI{5.000}{kg} & Locked aircraft mass used for lift, trim, climb, descent, and power calculations. \\
Maximum mass requirement & \SI{5.000}{kg} & Stage~1 mission-level requirement. \\
Payload mass & \SI{2.50}{lb} / \SI{1.134}{kg} & Delivery payload carried before drop-off. \\
Post-delivery flight mass & \SI{3.866}{kg} & Used for after-delivery performance reporting only. \\
Stage~1/2 component-estimated mass & \SI{2.446}{kg} & Frozen upstream propulsion, battery, and baseline system mass estimate. \\
Stage~3 component-estimated mass & \SI{2.491}{kg} & Stage~1/2 estimate plus the Stage~3 NGX250 foam tail estimate. \\
Remaining mass to gross target & \SI{2.509}{kg} & Payload, unmodeled structure, wiring, fasteners, covering, ballast, and contingency. \\
Main airfoil & DAE51 & Frozen Stage~1/2 wing workflow selection. \\
Tail airfoil & NACA0008 & Fixed NACA~0008 for both stabilizers. \\
Main wing span, chord, area & \SI{2.000}{m}, \SI{0.350}{m}, \SI{0.700}{m^2} & Frozen Stage~1/2 rectangular main wing. \\
Main wing aspect ratio & 5.714 & Derived from Stage~1/2 wing dimensions. \\
Main wing incidence & \SI{4.500}{deg} & Optimized by Stage~3 for cruise body-alpha compatibility. \\
Main wing washout & \SI{0.000}{deg} & Editable Stage~3 constraint; currently held fixed. \\
Propulsion layout & 10 \(\times\) \SI{5.50}{in}, \(P/D=0.40\), balanced & Frozen Stage~1/2 motor/propeller selection. \\
Stage~1/2 seed RPM, slow/cruise & \num{10062} / \num{9920} & Upstream RPMs retained for traceability; Stage~3 recomputes RPM using the active eCalc data. \\
Stage~3 RPM, slow/cruise & \num{8457} / \num{6976} & Final propeller operating points from the active Stage~3 propulsion model. \\
Power loss factors, wiring/avionics & 1.080 / 1.100 & Applied after eCalc/generic motor power to include battery harness, connector, distribution, BEC/regulator, and avionics wiring losses. \\
Propeller axial station & \SI{-0.0875}{m} & Relative to wing leading edge; negative is forward of the leading edge. \\
Fuselage nose and cross-section & \(x_n=\SI{-0.300}{m}\), \SI{170}{mm} \(\times\) \SI{130}{mm} & Nose is 300~mm ahead of the wing leading edge; cross-section from the current packaging assumption. \\
Flap span/chord/deflection & 0.650 semispan, 0.340 chord, \SI{40.0}{deg} & Frozen Stage~1/2 slotted-flap sizing. \\
Aileron span/chord & 0.280 semispan, 0.280 chord & Frozen Stage~1/2 aileron sizing. \\
Horizontal tail span and area & \SI{0.606}{m}, \SI{0.122}{m^2} & Stage~3 tail-volume and cruise-trim sizing. \\
Horizontal tail root/tip chord & \SI{0.202}{m} / \SI{0.202}{m} & Includes taper as an optimized/exported dimension. \\
Horizontal tail incidence & \SI{-4.747}{deg} & Cruise trim setting. \\
Elevator chord fraction & 0.220 & Sized for slow-flight pitch authority after stabilizer sizing. \\
Vertical tail height and area & \SI{0.377}{m}, \SI{0.055}{m^2} & Stage~3 directional-stability and control sizing. \\
Vertical tail root/tip chord & \SI{0.152}{m} / \SI{0.137}{m} & Includes taper as an optimized/exported dimension. \\
Vertical tail incidence & \SI{0.000}{deg} & Fixed yaw/incidence setting. \\
Rudder chord fraction & 0.353 & Sized for slow-flight yaw authority after stabilizer sizing. \\
Tail arm & \SI{1.000}{m} & Wing aerodynamic center to tail aerodynamic center. \\
CG location & \SI{0.0875}{m} / 25.00\% MAC & Enforced quarter-chord final CG target. \\
\bottomrule
\end{longtable}

\begin{longtable}{@{}p{0.25\linewidth}p{0.17\linewidth}p{0.18\linewidth}p{0.18\linewidth}p{0.14\linewidth}@{}}
\caption{Executive system-performance summary by flight condition.}\label{tab:exec-performance}\\
\toprule
Metric & Cruise & Slow flight & Approach/descent & Climb \\
\midrule
\endfirsthead
\toprule
Metric & Cruise & Slow flight & Approach/descent & Climb \\
\midrule
\endhead
Configuration & Clean & Flaps down & Flaps down & Clean \\
Freestream speed & \SI{10.00}{m/s} & \SI{4.00}{m/s} & \SI{5.00}{m/s} & \SI{10.00}{m/s} \\
Vertical rate target & -- & -- & \SI{6.00}{ft/s} sink & \SI{6.00}{ft/s} climb \\
Actual vertical rate & -- & -- & \SI{6.00}{ft/s} sink & \SI{6.00}{ft/s} climb \\
Flight-path angle & Level & Level & \SI{21.45}{deg} down & \SI{10.54}{deg} up \\
Free-stream / local \(C_L\) & 1.144 / 0.501 & 7.148 / 0.839 & 4.258 / 0.392 & 1.124 / 0.341 \\
Blown wing velocity & \SI{18.15}{m/s} & \SI{15.21}{m/s} & \SI{21.58}{m/s} & \SI{22.60}{m/s} \\
Effective dynamic pressure & \SI{139.8}{Pa} & \SI{83.5}{Pa} & \SI{166.2}{Pa} & \SI{201.9}{Pa} \\
Total drag & \SI{6.118}{N} & \SI{11.793}{N} & \SI{45.399}{N} & \SI{5.949}{N} \\
Required thrust & \SI{6.240}{N} & \SI{11.789}{N} & \SI{27.464}{N} & \SI{14.916}{N} \\
Throttle estimate & 34.8\% & 44.9\% & 100.0\% & 63.6\% \\
Motor/prop RPM & \num{6976} & \num{8457} & \num{12440} & \num{10134} \\
Electrical power & \SI{112.10}{W} & \SI{188.42}{W} & \SI{592.04}{W} & \SI{316.72}{W} \\
Aircraft electrical losses & \SI{8.49}{W} & \SI{14.14}{W} & Included in power & Included in power \\
Energy budgeted & \SI{37.37}{Wh} & \SI{4.71}{Wh} & Not in pack sizing & Not in pack sizing \\
Trim state & Body \(\alpha\) \SI{1.79}{deg}; \(i_h\) \SI{-4.75}{deg} & \(i_h\) \SI{-4.75}{deg}; flaps \SI{40}{deg} & Elevator \SI{15.25}{deg}; rudder \SI{0.00}{deg} & Elevator \SI{4.86}{deg} \\
\bottomrule
\end{longtable}

\begin{table}[H]
\centering
\small
\begin{tabular}{lcc}
\toprule
Quantity & Value & Calculation note \\
\midrule
Mass case & Payload on, \SI{5.000}{kg} & Takeoff is sized with payload still onboard. \\
Assumed climb speed & \SI{10.00}{m/s} & Same speed as the clean climb trim check. \\
Available takeoff RPM & \num{12440} & Maximum RPM in the active eCalc static thrust/power table. \\
Expected electrical power at max RPM & \SI{592.04}{W} & Includes Stage~3 wiring and avionics loss factors. \\
Expected 4S current at max RPM & \SI{40.0}{A} & \(I=P/V_{\mathrm{pack}}\), using \SI{14.8}{V} nominal. \\
Estimated max available thrust & \SI{23.47}{N} & Derived from the clean-climb thrust/throttle calibration. \\
Clean climb drag at \SI{10}{m/s} & \SI{5.95}{N} & Uses the current payload-on clean climb drag buildup. \\
Estimated excess thrust & \SI{17.52}{N} & \(T_{\max}-D_{\mathrm{climb}}\). \\
Estimated max climb angle & \SI{20.94}{deg} & \(\gamma_{\max}=\sin^{-1}\left((T_{\max}-D)/W\right)\). \\
Estimated max takeoff climb rate & \SI{3.57}{m/s} / \SI{11.72}{ft/s} & \(\dot{h}_{\max}=V\sin\gamma_{\max}=V(T_{\max}-D)/W\). \\
Configured climb check & \SI{1.83}{m/s} / \SI{6.00}{ft/s} & Current YAML climb requirement; this is below the max-RPM estimate. \\
\bottomrule
\end{tabular}
\caption{Payload-on maximum takeoff climb-rate estimate. This is a first-pass steady climb calculation at the selected climb speed; it does not include ground-roll acceleration, wheel drag, launch transients, gusts, or degraded battery voltage.}
\label{tab:stage3-takeoff-max-climb}
\end{table}

\begin{longtable}{@{}p{0.18\linewidth}p{0.12\linewidth}p{0.18\linewidth}p{0.18\linewidth}p{0.18\linewidth}p{0.12\linewidth}@{}}
\caption{Payload-on and post-delivery performance comparison. Battery capacity is sized from the payload-on case.}\label{tab:exec-payload-performance}\\
\toprule
Case & Mass & Cruise power/RPM/throttle & Slow power/RPM/throttle & Climb power/RPM/throttle & Cruise/Climb pitch margin \\
\midrule
\endfirsthead
\toprule
Case & Mass & Cruise power/RPM/throttle & Slow power/RPM/throttle & Climb power/RPM/throttle & Cruise/Climb pitch margin \\
\midrule
\endhead
Payload on & \SI{5.000}{kg} & \SI{112.1}{W} / \num{6976} / 34.8\% & \SI{188.4}{W} / \num{8457} / 44.9\% & \SI{316.7}{W} / \num{10134} / 63.6\% & 126\% / 233\% \\
Post delivery & \SI{3.866}{kg} & \SI{75.3}{W} / \num{5948} / 23.5\% & \SI{128.4}{W} / \num{7348} / 33.2\% & \SI{213.9}{W} / \num{8857} / 46.7\% & 157\% / 273\% \\
\bottomrule
\end{longtable}

\begin{longtable}{@{}p{0.31\linewidth}p{0.26\linewidth}p{0.35\linewidth}@{}}
\caption{Executive margins and PDR review notes.}\label{tab:exec-margins}\\
\toprule
Review item & Current value & Note \\
\midrule
\endfirsthead
\toprule
Review item & Current value & Note \\
\midrule
\endhead
Static margin & 0.237 MAC & Required range 0.100--0.270 MAC; target 0.130 MAC. \\
CG error from quarter chord & \SI{0.000000}{m} & Positive means aft of target; current target is 25.00\% MAC. \\
Horizontal tail volume & 0.500 & Constraint range 0.500--0.950. \\
Vertical tail volume & 0.039 & Constraint range 0.039--0.041. \\
Slow-flight lift margin & 86.4\% & Margin after the 1.20 stall-margin factor is applied. \\
Slow blown-panel local \(C_L\) margin & 0.764 & Positive means the blown panel is below its design local \(C_{L,\max}\). \\
Slow required blown-velocity margin & \SI{4.334}{m/s} & Actual minus required slipstream velocity at the slow-flight sizing point. \\
Main-wing pitch-up stall margin rule & 1.25x & Review assumption: operate at no more than \(1/1.25\) of local \(C_{L,\max}\) for clean emergency pitch-up margin. \\
Payload-on clean pitch margins & 126.5\% cruise / 232.7\% climb & Positive margin indicates available local \(C_L\) headroom beyond the review rule. \\
Post-delivery clean pitch margins & 156.6\% cruise / 273.5\% climb & Lower weight increases clean-wing pitch-up headroom. \\
Added drag target for slow flight & \SI{1.274}{N} & First-pass airbrake/spoiler drag requirement so blown-lift thrust does not accelerate the aircraft. \\
Pitch control margin & 212.5\% & Elevator authority margin at slow flight. \\
Yaw control margin & 0.0\% & Rudder authority margin at slow flight. \\
Approach RPM feasible flag & 0 & 1 means the requested thrust is inside the active RPM bracket. \\
Climb RPM feasible flag & 1 & 1 means the requested thrust is inside the active RPM bracket. \\
Battery capacity requirement & \SI{52.60}{Wh} / \num{3554}~mAh & 4S LiPo equivalent including 20\% reserve for \SI{20.0}{min} cruise plus \SI{1.5}{min} slow flight. \\
Battery mass estimate & \SI{0.292}{kg} & Uses Stage~1 pack specific energy of \SI{180}{Wh/kg}. \\
Design warnings & \texttt{slow\_flight\_requires\_powered\_lift;stage12\_equivalent\_clmax\_capped;approach\_throttle\_saturated} & Generated warning flags from the Stage~3 result row. \\
\bottomrule
\end{longtable}

\section*{Nomenclature}
\begin{longtable}{@{}ll@{}}
\toprule
Symbol & Meaning \\
\midrule
\(b,c,S\) & main-wing span, chord, and reference area \\
\(S_h,S_v\) & horizontal and vertical tail planform areas \\
\(b_h,b_v\) & horizontal tail span and vertical tail height \\
\(l_t\) & tail arm from wing aerodynamic center to tail aerodynamic center \\
\(V_h,V_v\) & horizontal and vertical tail volume coefficients \\
\(x_{cg},x_{np}\) & longitudinal center of gravity and neutral point \\
\(SM\) & static margin, \((x_{np}-x_{cg})/c\) \\
\(C_{L,\max}\) & freestream-reference maximum lift coefficient \\
\(C_{dA,f}\) & equivalent fuselage drag area \\
\(q=\frac{1}{2}\rho V^2\) & dynamic pressure \\
\bottomrule
\end{longtable}

\section{Scope and Frozen Inputs}
Stage~3 is not a re-optimization of the whole aircraft. The main wing and propulsion architecture are inherited from the Stage~1/2 workflow, and the propeller count is held fixed. The active stabilizer-planform design variables are the horizontal tail area, horizontal tail aspect ratio, horizontal tail taper ratio, vertical tail area, vertical tail aspect ratio, vertical tail taper ratio, and tail arm. Elevator and rudder chord fractions are not cruise-optimization variables; they are sized after planform selection from slow-speed pitch and yaw authority targets. The run is generated by \texttt{optimizer/stages/stage3\_aerosandbox.py}; editable Stage~3 constraints are read from \texttt{optimizer/config/stage3\_constraints.yaml}.

\begin{table}[H]
\centering
\begin{tabular}{llc}
\toprule
Quantity & Source & Value \\
\midrule
Gross flight mass & mission config & \SI{5.000}{kg} \\
Maximum built mass & mission config & \SI{5.000}{kg} \\
Main wing span \(b\) & Stage 1/2 fixed & \SI{2.000}{m} \\
Main wing chord \(c\) & Stage 1/2 fixed & \SI{0.350}{m} \\
Main wing area \(S\) & derived & \SI{0.700}{m^2} \\
Main airfoil & wing workflow & DAE51 \\
Tail airfoil & Stage 3 config & naca0008 \\
Propeller layout & Stage 1/2 fixed & 10 \(\times\) \SI{5.50}{in}, \(P/D=0.40\) \\
Propulsion model & Stage 3/eCalc & eCalc static table: outputs/ecalc\_prop\_analysis/x2302\_1500kv\_3s\_5p5x3p5\_3b/ecalc\_static\_partial\_load.csv \\
Power loss factors & Stage 3 config & wiring 1.080; avionics 1.100 \\
Fuselage nose station & Stage 3 config & \SI{-0.300}{m} from wing leading edge \\
Stage 3 slow-flight condition & Stage 3 config & \SI{4.00}{m/s} \\
Upstream Stage 1/2 low-speed condition & mission config & \SI{4.00}{m/s} \\
Cruise condition & mission config & \SI{10.00}{m/s} \\
Fuselage cross-section & Stage 3 config & \SI{170}{mm} \(\times\) \SI{130}{mm} \\
NGX250 foam density & material input & \SI{25.0}{kg/m^3} \\
\bottomrule
\end{tabular}
\caption{Frozen aircraft-level inputs and material assumptions used by Stage~3.}
\label{tab:stage3-fixed-inputs}
\end{table}

\section{How to Run and Reproduce This Stage}
Stage~3 is intended to be run from the repository root after Stage~1 and Stage~2 have produced their CSV outputs. The command is
\begin{verbatim}
python3 -m optimizer.stages.stage3_aerosandbox
\end{verbatim}
A successful run writes the queue CSV, the full results CSV, the top-design CSV, the readable Markdown summary, this LaTeX engineering report, and all geometry and performance figures. The recommended post-run validation command is
\begin{verbatim}
python3 -m optimizer.stages.validate_outputs
\end{verbatim}
The editable user-facing inputs are stored in \texttt{optimizer/config/stage3\_constraints.yaml}. Values in that YAML file are loaded into the Stage3SizingConfig dataclass. Unknown keys are ignored, so comments and temporary notes can be added without breaking the run.

\subsection{Data flow}
\begin{enumerate}[leftmargin=*]
\item Stage~3 reads \texttt{outputs/stage2\_prop\_span\_report.csv} and filters it to the frozen propulsion layout selected by the Stage~1/2 wing workflow.
\item It reads \texttt{outputs/stage1\_pareto\_front.csv} to recover the matching mass, drag, propeller, battery, and RPM context for that frozen layout.
\item It reads the wing workflow summary to recover the selected airfoil, flap geometry, aileron geometry, propeller spanwise positions, and the Stage~1/2 \(C_{L,\max}\) curves.
\item If enabled in the YAML constraints and the frozen propulsion layout matches, it reads the eCalc static thrust/power table and dynamic design-point \(C_T/C_P\) CSV files to replace the generic propeller surrogate used in early screening.
\item After the propulsion model returns motor/ESC electric power, Stage~3 applies the editable wiring and avionics loss factors from the YAML file. The reported condition powers and battery-capacity estimate therefore include harness, connector, distribution, BEC/regulator, and avionics wiring allowances.
\item It optimizes only the tail planform and tail arm. The main wing, propeller count, propeller size, propeller pitch ratio, flap sizing, aileron sizing, and mass target stay fixed.
\item It sizes the elevator and rudder chord fractions after the stabilizer planform is selected, because those moving surfaces are judged mainly by slow-flight control authority rather than cruise drag.
\item It renders geometry figures, writes sweep plots, writes drag-component plots, and exports the final design rows to the Stage~3 CSV outputs.
\end{enumerate}

\subsection{Primary output files}
\begin{itemize}[leftmargin=*]
\item \texttt{outputs/stage3\_readable\_results.md}: concise engineering summary for humans.
\item \texttt{outputs/stage3\_aerosandbox\_results.csv}: full machine-readable results.
\item \texttt{outputs/stage3\_aerosandbox\_top\_designs.csv}: first ranked designs.
\item \texttt{outputs/stage3\_visuals/*drag\_power\_sweeps.png}: drag and power curves.
\item \texttt{outputs/stage3\_visuals/*drag\_components.png}: drag-buildup plot.
\item \texttt{outputs/stage3\_visuals/*performance\_sweep.csv}: underlying sweep data.
\item \texttt{outputs/stage3\_visuals/*mesh.npz}: AeroSandbox mesh points and faces.
\end{itemize}

\section{Modeling Method}
\subsection{Tail geometry parameterization}
Both stabilizers are modeled as tapered trapezoids. The optimizer chooses area, aspect ratio, taper ratio, and tail arm inside the user-specified tail-volume ranges. The span and root/tip chords are then computed from
\begin{align}
b_h &= \sqrt{S_h AR_h}, & c_{r,h} &= \frac{2S_h}{b_h(1+\lambda_h)}, & c_{t,h} &= \lambda_h c_{r,h}, \\
b_v &= \sqrt{S_v AR_v}, & c_{r,v} &= \frac{2S_v}{b_v(1+\lambda_v)}, & c_{t,v} &= \lambda_v c_{r,v}.
\end{align}
The horizontal and vertical tail volume coefficients are evaluated as
\begin{equation}
V_h = \frac{S_h l_t}{S c}, \qquad V_v = \frac{S_v l_t}{S b}.
\end{equation}
For this run, the requested horizontal tail-volume range is 0.500 to 0.950, and the requested vertical tail-volume range is 0.039 to 0.041.

\subsection{Optimization objective and penalties}
The optimizer uses a scalar score rather than a hard single-objective aerodynamic minimum. This is deliberate: a tail with very low drag but poor static margin, unreasonable incidence, excessive mass, or inadequate slow-flight lift margin is not a useful aircraft design. The score has normalized cruise-power, slow-flight-power, tail-mass, and static-margin terms:
\begin{equation}
J = w_c \frac{P_c}{P_{c,ref}} + w_s \frac{P_s}{P_{s,ref}} + w_m \frac{m_t}{m_{max}} + w_{SM}\left(\frac{SM-SM_{target}}{0.08}\right)^2 + J_{penalty}.
\end{equation}
The current weights are \(w_c=0.68\), \(w_s=0.06\), \(w_m=0.16\), and \(w_{SM}=0.10\). Cruise power is weighted most heavily because the stabilizer planform is supposed to be cruise-oriented. Slow-flight control authority is not ignored; it is handled by the post-optimization elevator and rudder sizing step.
Penalty terms are applied when a candidate falls outside the configured tail-volume ranges, static-margin range, mass budget, lift-margin requirement, trim-incidence limit, root-chord bounds, or propeller RPM feasibility. This keeps the optimizer from using an attractive but non-buildable corner of the design space.

\subsection{Slow-flight elevator and rudder sizing}
After the cruise-oriented stabilizer planform is selected, Stage~3 sizes the elevator and rudder chord fractions against slow-flight control-authority targets. The authority model uses finite-wing lift slope, tail volume, maximum deflection, slow-flight tail dynamic-pressure ratio, and a bounded control-surface effectiveness model:
\begin{align}
C_{m,\delta_e}^{\max} &\approx \eta_{t,\mathrm{slow}} a_h V_h \tau_e \delta_{e,\max}, \\
C_{n,\delta_r}^{\max} &\approx \eta_{t,\mathrm{slow}} a_v V_v \tau_r \delta_{r,\max}.
\end{align}
This keeps the stabilizer dimensions focused on cruise performance while still ensuring that the moving surfaces retain low-speed control authority.

\subsection{Foam mass accounting}
The NGX250 foam mass is included during the tail sizing loop. For a trapezoidal foam surface with airfoil-area coefficient \(k_A\), span \(b_s\), root chord \(c_r\), and tip chord \(c_t\), Stage~3 uses
\begin{equation}
m_{\mathrm{foam}} = \rho_{\mathrm{foam}} k_A b_s \frac{c_r^2 + c_r c_t + c_t^2}{3}.
\end{equation}
The tail foam mass is added to the selected Stage~1/2 built-mass estimate before checking the mass margin and recomputing the longitudinal center of gravity.

\subsection{Trim, static margin, and drag buildup}
Cruise trim is solved with AeroSandbox VortexLatticeMethod. A proxy neutral point is also evaluated during optimization using finite-wing lift slopes, tail dynamic-pressure ratio, downwash gradient, and tail arm. The static margin constraint is
\begin{equation}
SM = \frac{x_{np}-x_{cg}}{c}.
\end{equation}
Cruise drag is assembled from main-wing profile and induced drag, horizontal-tail trim drag, vertical-tail profile drag, and an equivalent fuselage parasite increment:
\begin{equation}
D_{\mathrm{cruise}} = qS(C_{D,p}+C_L^2/\pi e AR) + qS_h C_{D,h} + qS_v C_{D,v} + q C_{dA,f}.
\end{equation}
For this run, \(C_{dA,f}=\SI{0.011}{m^2}\), \(q_{\mathrm{low}}=\SI{9.800}{Pa}\), and \(q_{\mathrm{cruise}}=\SI{61.250}{Pa}\).

\subsection{Slow-flight high-lift treatment}
Slow flight is always evaluated with slotted flaps deployed. Stage~3 no longer invents a large physical \(C_{L,\max}\) for this mode. It reads the selected Stage~1/2 equivalent high-lift outputs for traceability, and it reads the Stage~1/2 NeuralFoil flap-sweep section polars for the physical local clean and flapped \(C_{L,\max}\) limits. The YAML CLmax values are only missing-file fallbacks. Propwash increases local dynamic pressure rather than airfoil CLmax. The loaded sources are
\begin{itemize}[leftmargin=*]
\item \texttt{outputs/wing\_workflow/dae51/control\_surface\_sizing/rank01\_n10\_d5p5\_p2p2\_balanced\_b3/dae51/control\_surface\_summary.csv}
\item \texttt{outputs/wing\_workflow/dae51/control\_surface\_sizing/rank01\_n10\_d5p5\_p2p2\_balanced\_b3/dae51/total\_cl\_curve.csv}
\item \texttt{outputs/wing\_workflow/dae51/control\_surface\_sizing/rank01\_n10\_d5p5\_p2p2\_balanced\_b3/dae51/flap\_sweep.csv}
\end{itemize}
The Stage~1/2 section-polar values are \(C_{L,\max,clean}\)=1.418/1.433 (unblown/blown) and \(C_{L,\max,flap}\)=2.152/2.171. The area-weighted physical blown-panel limit used by Stage~3 is 1.876.
\begin{table}[H]
\centering
\begin{tabular}{lcc}
\toprule
Condition & Raw equivalent \(C_L\) & Physical cap used \\
\midrule
No flap & 1.418 & 1.418 \\
Flap only & 1.876 & 1.876 \\
Clean blown & 5.545 & 1.433 \\
Flap-down blown & 5.828 & 1.876 \\
\bottomrule
\end{tabular}
\caption{Raw Stage~1/2 equivalent high-lift coefficients and the physical local CLmax caps used by Stage~3.}
\label{tab:stage3-clmax}
\end{table}
The slow-flight sizing coefficients are divided by the stall margin factor 1.20. Thus the design flap-only coefficient is 1.564. The local physical blown-panel design coefficient is 1.563, while the freestream-referenced blown equivalent design coefficient is 13.321. The current slow-flight blown-panel local CL is 0.799.

\subsection{Slow-flight force balance and added drag}
The blown-lift calculation can require more propeller thrust than the natural aircraft drag at the low-speed condition. If that excess thrust is not cancelled, the aircraft would accelerate and would no longer remain at the intended slow-flight speed. Stage~3 therefore reports the drag that must be added by future drag devices.
\begin{align}
D_{\mathrm{natural,slow}} &= D_{\mathrm{Stage1/2,flap}} + D_h + D_v + D_f, \\
T_{\mathrm{blown}} &= \max(T_{D,\mathrm{slow}}, T_{V_{\mathrm{eff}}}), \\
\Delta D_{\mathrm{add}} &= \max(0, T_{\mathrm{blown}} - D_{\mathrm{natural,slow}}), \\
D_{\mathrm{steady,slow}} &= D_{\mathrm{natural,slow}} + \Delta D_{\mathrm{add}}.
\end{align}
For this design, the natural slow-flight drag is \SI{10.519}{N}, the propeller thrust needed by the blown-lift condition is \SI{11.793}{N}, and the added-drag target is \SI{1.274}{N}. This number should be treated as the first sizing target for airbrakes, spoilerons, deployable plates, or other high-drag devices.

\subsection{Performance sweeps}
The drag and power sweep figure is not part of the optimizer objective; it is a diagnostic output for the final ranked design. The clean velocity sweep solves the angle of attack needed to support aircraft weight and omits points outside the clean polar lift range. The flaps-down sweeps use the same split blown/unblown component model as the Stage~3 slow-flight sizing calculation, with profile drag, induced drag, tail drag, fuselage drag, and propeller power evaluated from that model.
The power curves no longer plot saturated max-RPM points as valid requirements. If the requested trim, drag balance, or blown-lift slipstream condition cannot be met within the configured RPM range, the point is left blank in the plot and marked with zero feasibility in the sweep CSV.

\section{Assumption Justification}
The assumptions below are intentionally exposed in the YAML file when they are likely to change during design iteration. The current values are first-pass engineering choices, not immutable truth.
\begin{longtable}{@{}p{0.24\linewidth}p{0.18\linewidth}p{0.49\linewidth}@{}}
\toprule
Assumption & Current value & Justification \\
\midrule
Tail airfoil & naca0008 & Fixed for both stabilizers. A symmetric 8\% section gives near-zero airfoil pitching moment, lower profile thickness than NACA~0012, and enough foam thickness for a small tail. \\
Foam density & \SI{25.0}{kg/m^3} & NGX250 pink foam is the selected construction material, so tail volume is converted directly into added mass. \\
Horizontal tail-volume range & 0.500--0.950 & Conventional preliminary sizing range for a small stable aircraft, widened enough for the high-lift/slow-flight requirement. \\
Vertical tail-volume range & 0.039--0.041 & Keeps the fin large enough for yaw control without allowing unnecessary cruise drag and foam mass. \\
Static-margin target & 0.130 MAC & A moderate stability target. The lower bound avoids neutral/unstable pitch behavior; the upper bound avoids an overly nose-heavy aircraft with excess trim drag. \\
Cruise tail dynamic-pressure ratio & 0.90 & The tail is assumed to see slightly less than freestream dynamic pressure due to wing wake, fuselage interference, and downwash. \\
Slow tail dynamic-pressure ratio & 0.85 & Slow-flight control sizing is made more conservative because flap wake, high angle of attack, and low Reynolds number can reduce tail effectiveness. \\
Downwash gradient & 0.35 & A first-pass low-aspect-ratio wing estimate for how much wing angle of attack becomes downwash at the tail. \\
Fuselage drag area & \SI{0.011}{m^2} & Equivalent drag area for the \SI{170}{mm} by \SI{130}{mm} fuselage cross-section with a modest bluff-body/protuberance allowance. \\
Tail zero-lift drag coefficient & 0.012 & Conservative small-tail, low-Reynolds-number profile drag estimate including practical surface roughness and hinge gaps. \\
Stall margin factor & 1.20 & Divides inherited Stage~1/2 \(C_{L,\max}\) by 1.20 so slow-flight sizing keeps 20\% lift headroom for gusts, manufacturing errors, surface waviness, and blown-lift uncertainty. \\
Minimum slow-lift warning margin & 0.10 & Adds a warning if less than 10\% lift margin remains after the stall margin has already been applied. \\
Elevator/rudder max deflection & \SI{25}{deg} / \SI{25}{deg} & Avoids assuming extreme deflections where hinge moments, separation, and servo loads become questionable. \\
Control effectiveness exponent & 0.70 & Represents diminishing returns as control-surface chord fraction grows; authority is not assumed to scale linearly forever. \\
Wiring power loss factor & 1.080 & Multiplies propulsion electric power after the eCalc/generic motor model to cover battery harness, connector, and distribution losses that are not part of the aerodynamic thrust requirement. \\
Avionics power loss factor & 1.100 & Multiplies the fixed avionics load to include BEC/regulator inefficiency and wiring loss in the avionics bus. \\
Approach target & \SI{6.0}{ft/s}, \SI{5.0}{m/s} & Editable first-pass descent condition. If the target sink rate exceeds what the selected approach speed can physically provide, Stage~3 clips it and reports that the target was not met. \\
Battery sizing segments & \SI{20.0}{min} cruise, \SI{1.5}{min} slow & Converts the final Stage~3 operating powers into a pack-capacity requirement. The reserve fraction and specific energy are inherited from the Stage~1 mission assumptions so the battery estimate remains tied to the same aircraft-level mass model. \\
\bottomrule
\end{longtable}

\section{Recommended Stage 3 Geometry}
\begin{table}[H]
\centering
\begin{tabular}{lcc}
\toprule
Quantity & Symbol & Value \\
\midrule
Main wing span & \(b\) & \SI{2.000}{m} \\
Main wing chord & \(c\) & \SI{0.350}{m} \\
Main wing incidence & \(i_w\) & \SI{4.500}{deg} \\
Cruise body alpha proxy & \(\alpha_b\) & \SI{1.790}{deg} \\
Cruise section alpha proxy & \(\alpha_\mathrm{sec}\) & \SI{6.290}{deg} \\
Propeller axial position & \(x_p\) & \SI{-0.0875}{m} \\
Fuselage nose station & \(x_n\) & \SI{-0.3000}{m} \\
Horizontal tail span & \(b_h\) & \SI{0.606}{m} \\
Horizontal tail root chord & \(c_{r,h}\) & \SI{0.202}{m} \\
Horizontal tail tip chord & \(c_{t,h}\) & \SI{0.202}{m} \\
Horizontal tail taper & \(\lambda_h\) & 1.000 \\
Horizontal tail incidence & \(i_h\) & \SI{-4.747}{deg} \\
Elevator chord fraction & \(c_e/c_h\) & 0.220 \\
Vertical tail height & \(b_v\) & \SI{0.377}{m} \\
Vertical tail root chord & \(c_{r,v}\) & \SI{0.152}{m} \\
Vertical tail tip chord & \(c_{t,v}\) & \SI{0.137}{m} \\
Vertical tail taper & \(\lambda_v\) & 0.902 \\
Vertical tail incidence & \(i_v\) & \SI{0.000}{deg} \\
Rudder chord fraction & \(c_r/c_v\) & 0.353 \\
Tail arm & \(l_t\) & \SI{1.000}{m} \\
\bottomrule
\end{tabular}
\caption{Final Stage~3 aircraft dimensions.}
\label{tab:stage3-dimensions}
\end{table}

\begin{figure}[H]
\centering
\includegraphics[width=0.72\linewidth]{rank01_n10_d5p5_pd0p40_balanced_top_view.png}
\caption{Stage~3 top-view fixed-wing tail layout. The dashed circles show the actual frozen Stage~2 propeller count and spanwise positions; the axial position is read from the Stage~3 YAML constraint file rather than hard-coded in the renderer.}
\label{fig:stage3-top-view}
\end{figure}

\begin{figure}[H]
\centering
\includegraphics[width=0.92\linewidth]{rank01_n10_d5p5_pd0p40_balanced_three_view.png}
\caption{AeroSandbox-generated three-view rendering of the final Stage~3 geometry.}
\label{fig:stage3-three-view}
\end{figure}

\begin{figure}[H]
\centering
\includegraphics[width=0.82\linewidth]{rank01_n10_d5p5_pd0p40_balanced_3d_wireframe.png}
\caption{Three-dimensional wireframe rendering used as a geometry sanity check for the final Stage~3 aircraft.}
\label{fig:stage3-3d}
\end{figure}

\section{Mass, Stability, and Performance Results}
\begin{table}[H]
\centering
\begin{tabular}{lcc}
\toprule
Metric & Value & Comment \\
\midrule
Stage~1/2 built mass & \SI{2.446}{kg} & inherited baseline \\
Horizontal tail foam mass & \SI{0.034}{kg} & NGX250, geometry-derived \\
Vertical tail foam mass & \SI{0.011}{kg} & NGX250, geometry-derived \\
Total tail foam mass & \SI{0.045}{kg} & added during sizing \\
Gross flight mass used in aerodynamics & \SI{5.000}{kg} & locked Stage~1/2 flight mass \\
Component-estimated mass before ballast/payload & \SI{2.491}{kg} & baseline plus tail foam \\
Remaining mass to gross target & \SI{2.509}{kg} & payload, unmodeled structure, wiring, fasteners, ballast, contingency \\
Mass margin & \SI{2.509}{kg} & relative to \SI{5.0}{kg} limit \\
CG location & \SI{0.088}{m} & from wing leading edge \\
CG fraction MAC & 25.00\% & based on rectangular MAC \\
CG target & 25.00\% MAC & quarter-chord requirement \\
CG error & \SI{0.000000}{m} & actual minus target \\
Required pre-tail baseline CG & 19.77\% MAC & mass placement before tail foam \\
Neutral point & \SI{0.170}{m} & proxy stability model \\
Static margin & 0.237 MAC & target 0.237 MAC achieved \\
Horizontal tail volume & 0.500 & range 0.500--0.950 \\
Vertical tail volume & 0.039 & range 0.039--0.041 \\
VLM \(C_{m,\alpha}\) & -0.01750/deg & negative is statically stable \\
\bottomrule
\end{tabular}
\caption{Mass and static-stability closure for the selected Stage~3 geometry.}
\label{tab:stage3-mass-stability}
\end{table}

\begin{table}[H]
\centering
\begin{tabular}{lccc}
\toprule
Control axis & Target & Available & Margin \\
\midrule
Pitch/elevator & \(C_m=0.180\) & \(C_m=0.563\) & 212.5\% \\
Yaw/rudder & \(C_n=0.055\) & \(C_n=0.055\) & 0.0\% \\
\bottomrule
\end{tabular}
\caption{Slow-flight control-surface sizing results after cruise-oriented stabilizer planform selection.}
\label{tab:stage3-control-authority}
\end{table}

\begin{table}[H]
\centering
\begin{tabular}{lcc}
\toprule
Metric & Low speed & Cruise \\
\midrule
Freestream speed & \SI{4.00}{m/s} & \SI{10.00}{m/s} \\
Total drag & \SI{11.793}{N} & \SI{6.118}{N} \\
Natural airframe drag & \SI{10.519}{N} & -- \\
Slow wing induced, used & \SI{3.618}{N} & -- \\
Slow wing induced, local blown-panel & \SI{3.618}{N} & -- \\
Slow wing induced, freestream-equivalent & \SI{29.621}{N} & -- \\
Added drag required & \SI{1.274}{N} & -- \\
Slow drag minus cruise drag & \SI{5.675}{N} & -- \\
Fuselage drag increment & \SI{0.108}{N} & \SI{0.674}{N} \\
Electrical power & \SI{188.42}{W} & \SI{112.10}{W} \\
Freestream-reference \(C_L\) & 7.148 & 1.144 \\
Pressure-weighted local \(C_L\) & 0.839 & 0.501 \\
Unblown/blown dynamic pressure & \SI{9.8}{Pa} / \SI{141.7}{Pa} & \SI{61.3}{Pa} / \SI{201.8}{Pa} \\
Effective dynamic pressure & \SI{83.5}{Pa} & \SI{139.8}{Pa} \\
Required thrust & \SI{11.789}{N} & \SI{6.240}{N} \\
Throttle estimate & 44.9\% & 34.8\% \\
RPM & \num{8457} & \num{6976} \\
\(C_T/C_P\) & 0.1269 / 0.0706 & 0.0987 / 0.0706 \\
Required/actual blown velocity & \SI{10.874}{m/s} / \SI{15.208}{m/s} & -- \\
Cruise \(C_L/C_D/L/D\) & -- & 1.144 / 0.1427 / 8.02 \\
\bottomrule
\end{tabular}
\caption{Low-speed and cruise performance after Stage~3 tail sizing.}
\label{tab:stage3-performance}
\end{table}

\begin{table}[H]
\centering
\small
\begin{tabular}{lrrrrrrr}
\toprule
Condition & Speed [m/s] & RPM & Thrust [N] & Throttle [\%] & \(C_L\) free/local & Drag [N] & Power [W] \\
\midrule
Clean cruise & 10.00 & \num{6976} & 6.240 & 34.8 & 1.144 / 0.501 & 6.118 & 112.10 \\
Flaps-down slow flight & 4.00 & \num{8457} & 11.789 & 44.9 & 7.148 / 0.839 & 11.793 & 188.42 \\
Flaps-down approach/descent & 5.00 & \num{12440} & 27.464 & 100.0 & 4.258 / 0.392 & 45.399 & 592.04 \\
Clean climb & 10.00 & \num{10134} & 14.916 & 63.6 & 1.124 / 0.341 & 5.949 & 316.72 \\
\bottomrule
\end{tabular}
\caption{Flight-condition operating values. The free-stream \(C_L\) uses only \(q_\infty\); the local/effective \(C_L\) includes the pressure-weighted blown and unblown wing sections where that model is available.}
\label{tab:stage3-cl-rpm}
\end{table}

\begin{table}[H]
\centering
\begin{tabular}{lrrrr}
\toprule
Condition & Blown velocity [m/s] & \(q_\infty\) [Pa] & \(q_\mathrm{blown}\) [Pa] & \(q_\mathrm{eff}\) [Pa] \\
\midrule
Clean cruise & 18.152 & 61.3 & 201.8 & 139.8 \\
Flaps-down slow flight & 15.208 & 9.8 & 141.7 & 83.5 \\
Flaps-down approach/descent & 21.580 & 15.3 & 285.2 & 166.2 \\
Clean climb & 22.604 & 61.3 & 312.9 & 201.9 \\
\bottomrule
\end{tabular}
\caption{Propwash dynamic-pressure diagnostics used to translate free-stream lift demand into local wing-section lift demand.}
\label{tab:stage3-propwash-q}
\end{table}

\begin{table}[H]
\centering
\begin{tabular}{lrrr}
\toprule
Segment & Duration [min] & Power [W] & Energy [Wh] \\
\midrule
Cruise & 20.0 & 112.10 & 37.37 \\
Slow flight & 1.5 & 188.42 & 4.71 \\
Usable subtotal & 21.5 & -- & 42.08 \\
Pack capacity with 20\% reserve & -- & -- & 52.60 \\
Equivalent 4S LiPo capacity & -- & -- & 3554 mAh at \SI{14.8}{V} \\
\bottomrule
\end{tabular}
\caption{Battery capacity estimate for \SI{20.0}{min} of cruise and \SI{1.5}{min} of slow flight. At \SI{180}{Wh/kg}, the reserve-adjusted pack mass estimate is \SI{0.292}{kg}.}
\label{tab:stage3-battery-capacity}
\end{table}

\begin{table}[H]
\centering
\begin{tabular}{lrrrrrrr}
\toprule
Regime & RPM & Time [min] & Total P [W] & Total I [A] & Energy [Wh] & Used [mAh] & \% of sized 4S pack \\
\midrule
Takeoff & \num{10134} & 0.75 & 317 & 21.4 & 3.96 & 268 & 9.41\% \\
Loiter pre-drop & \num{6976} & 10.00 & 112 & 7.6 & 18.68 & 1262 & 44.40\% \\
Drop / low-speed & \num{8457} & 1.00 & 188 & 12.7 & 3.14 & 212 & 7.46\% \\
Loiter post-drop & \num{5948} & 5.00 & 75 & 5.1 & 6.28 & 424 & 14.92\% \\
Landing & \num{7348} & 0.75 & 128 & 8.7 & 1.60 & 108 & 3.81\% \\
\textbf{Total} & -- & \textbf{17.50} & -- & -- & \textbf{33.66} & \textbf{2275} & \textbf{80.00\%} \\
\bottomrule
\end{tabular}
\caption{Mission-regime energy estimate in the review-table format. Takeoff uses the payload-on climb condition; landing uses the post-delivery flaps-down slow/flare condition. Current is computed at nominal pack voltage.}
\label{tab:stage3-mission-regime-energy}
\end{table}
The mission-regime table sizes to a \textbf{\num{2843}~mAh} 4S LiPo at \SI{14.8}{V} nominal, including 20\% reserve. The usable energy before reserve is \SI{33.66}{Wh}, and the reserve-adjusted capacity is \SI{42.08}{Wh}.

\begin{table}[H]
\centering
\begin{tabular}{lccc}
\toprule
Aero surface & Cruise Reynolds number & Slow-flight Reynolds number & Characteristic chord \\
\midrule
Main wing & \num{236878} & \num{94751} & main wing chord \\
Horizontal stabilizer & \num{136747} & \num{54699} & H-tail MAC \\
Vertical stabilizer & \num{98166} & \num{39266} & V-tail MAC \\
\bottomrule
\end{tabular}
\caption{Reynolds numbers for the main lifting surfaces. Stage~3 uses \(Re=\rho V c/\mu\), with the tail characteristic chord equal to each trapezoidal mean aerodynamic chord.}
\label{tab:stage3-reynolds}
\end{table}

\begin{table}[H]
\centering
\begin{tabular}{lcc}
\toprule
Drag component & Clean cruise [N] & Flaps-down slow [N] \\
\midrule
Main wing profile & 0.528 & 6.735 \\
Main wing induced, used & 4.739 & 3.618 \\
Horizontal tail & 0.136 & 0.042 \\
Vertical tail & 0.040 & 0.016 \\
Fuselage & 0.674 & 0.108 \\
Added drag / airbrakes & 0.000 & 1.274 \\
Total & 6.118 & 11.793 \\
\bottomrule
\end{tabular}
\caption{Drag buildup for the clean cruise and flaps-down slow-flight sizing points. The slow-flight induced value is the physical split-panel sum; the freestream-equivalent value is retained only as a diagnostic check.}
\label{tab:stage3-drag-components}
\end{table}

\begin{figure}[H]
\centering
\includegraphics[width=0.92\linewidth]{rank01_n10_d5p5_pd0p40_balanced_drag_power_sweeps.png}
\caption{Total drag and electrical power required versus velocity and angle of attack for clean cruise and flaps-down slow-flight configurations.}
\label{fig:stage3-drag-power-sweeps}
\end{figure}

\begin{figure}[H]
\centering
\includegraphics[width=0.88\linewidth]{rank01_n10_d5p5_pd0p40_balanced_drag_components.png}
\caption{Drag component comparison for the final Stage~3 configuration.}
\label{fig:stage3-drag-components}
\end{figure}

\section{Approach Trim Estimate}
The approach calculation assumes a steady flaps-down descent at the user-specified target flight-path angle. The weight component along the flight path reduces the thrust required according to \(T = D - W\sin\gamma\), clipped at zero thrust. Elevator trim is estimated from the flapped pitching moment, the CG/aerodynamic-center offset, horizontal tail volume, slow-flight tail dynamic-pressure ratio, and the editable trim-effectiveness factor in the YAML constraints.

\begin{table}[H]
\centering
\begin{tabular}{lcc}
\toprule
Quantity & Symbol & Value \\
\midrule
Target sink rate & -- & \SI{6.00}{ft/s} \\
Actual sink rate & -- & \SI{1.829}{m/s} / \SI{6.00}{ft/s} \\
Sink-rate target met & -- & 1 \\
Approach angle below horizontal & \(\gamma\) & \SI{21.45}{deg} \\
Approach speed & \(V_A\) & \SI{5.00}{m/s} \\
Flap deflection & \(\delta_f\) & \SI{40.0}{deg} \\
Angle of attack & \(\alpha_A\) & \SI{13.06}{deg} \\
Required lift coefficient & \(C_L\) & 4.258 \\
Approach drag & \(D_A\) & \SI{45.399}{N} \\
Required thrust & \(T_A\) & \SI{27.464}{N} \\
Electrical power & \(P_A\) & \SI{592.04}{W} \\
Throttle estimate & -- & 100.0\% \\
Propulsion feasible & -- & 0 \\
RPM & -- & \num{12440} \\
Elevator trim & \(\delta_e\) & \SI{15.25}{deg} \\
Rudder trim & \(\delta_r\) & \SI{0.00}{deg} \\
Max descent rate limit & -- & \SI{6.00}{ft/s} \\
\bottomrule
\end{tabular}
\caption{Estimated trims and throttle setting for the configured approach angle.}
\label{tab:stage3-approach}
\end{table}

\section{Trim and Aerodynamic Verification}
The final geometry is checked with AeroSandbox VLM at cruise. The trim solver adjusts aircraft angle of attack and horizontal-tail incidence to close lift and pitching moment. The final residuals are
\begin{equation}
\Delta L = \SI{-0.000940}{N}, \qquad \Delta C_M = -0.000000.
\end{equation}
The airfoil-polar cruise body-angle proxy is \SI{1.790}{deg}, with a wing-section alpha proxy of \SI{6.290}{deg}. The VLM trim-alpha diagnostic is \SI{7.691}{deg}, and the horizontal-tail incidence is \SI{-4.747}{deg}. The selected result has warning status \texttt{slow\_flight\_requires\_powered\_lift;stage12\_equivalent\_clmax\_capped;approach\_throttle\_saturated}.

\begin{figure}[H]
\centering
\includegraphics[width=0.82\linewidth]{rank01_n10_d5p5_pd0p40_balanced_polars.png}
\caption{Section polar diagnostic used by Stage~3 for cruise and low-speed airfoil checks. The whole-wing slow-flight \(C_{L,\max}\) values used for sizing are inherited from the Stage~1/2 high-lift workflow rather than from this diagnostic plot alone.}
\label{fig:stage3-polars}
\end{figure}

\section{Engineering Interpretation and Open Items}
\begin{enumerate}[leftmargin=*]
\item The tail geometry meets the active tail-volume and static-margin targets while adding only \SI{0.045}{kg} of NGX250 foam mass.
\item The propeller drawing now uses the Stage~2 propeller centers and the Stage~3 YAML axial-location parameter. This makes the layout plot a geometry check rather than a schematic with a hard-coded propeller station.
\item The fuselage model has been updated to a \SI{170}{mm} by \SI{130}{mm} cross-section. The equivalent drag area remains a concept-level parasite-drag surrogate and should eventually be replaced by a measured or CFD-derived fuselage drag polar.
\item The slow-flight lift margin is large because Stage~3 consumes the Stage~1/2 flap-down blown \(C_{L,\max}\) directly. This is consistent with the current user requirement, but it also means the high-lift model remains one of the dominant validation risks.
\item The next fidelity step should be a coupled tail-plus-flap trim case at the low-speed condition, because the current high-lift sizing focuses on lift availability and propeller power rather than full low-speed moment trim with the flap deployed.
\end{enumerate}

\section{Generated Artifacts}
\begin{itemize}[leftmargin=*]
\item Machine-readable results: \texttt{../outputs/stage3\_aerosandbox\_results.csv}
\item Human-readable Markdown summary: \texttt{../outputs/stage3\_readable\_results.md}
\item Technical Markdown report: \texttt{../outputs/stage3\_aerosandbox\_report.md}
\item Top-view geometry plot: \texttt{rank01\_n10\_d5p5\_pd0p40\_balanced\_top\_view.png}
\item Three-view geometry plot: \texttt{rank01\_n10\_d5p5\_pd0p40\_balanced\_three\_view.png}
\item 3D wireframe rendering: \texttt{rank01\_n10\_d5p5\_pd0p40\_balanced\_3d\_wireframe.png}
\item Drag and power sweep plot: \texttt{rank01\_n10\_d5p5\_pd0p40\_balanced\_drag\_power\_sweeps.png}
\item Drag component plot: \texttt{rank01\_n10\_d5p5\_pd0p40\_balanced\_drag\_components.png}
\item Performance sweep CSV: \texttt{../outputs/stage3\_visuals/rank01\_n10\_d5p5\_pd0p40\_balanced\_performance\_sweep.csv}
\item Mesh file: \texttt{../outputs/stage3\_visuals/rank01\_n10\_d5p5\_pd0p40\_balanced\_mesh.npz}
\end{itemize}

\end{document}
